\documentclass[10pt]{article}
\usepackage[T1]{fontenc}
\usepackage[utf8]{inputenc}
\usepackage[french]{babel}
\usepackage{amsmath,amssymb,mathrsfs,mathpazo}
\usepackage{geometry}
\geometry{a4paper,margin=2cm}
\usepackage{graphicx,array,tabularx,tabularray,colortbl,multirow,multicol,booktabs,enumitem,pifont,ulem,fancybox,diagbox,cancel}
\usepackage{pstricks,pstricks-add,pst-plot,pst-tree,pst-node,pst-text,pst-func,pst-math,pst-bspline,pst-3dplot,pst-eucl,pst-solides3d}
\let\pscalculate\undefined
\usepackage{pst-bezier}
\usepackage{tikz,pgfplots,tkz-tab}
\usetikzlibrary{arrows,arrows.meta,patterns,calc,positioning,trees,decorations.markings,shapes.geometric,shapes.misc,shapes.arrows,intersections,decorations.pathmorphing,decorations.pathreplacing,backgrounds,shadows,babel}
\pgfplotsset{compat=1.18}
\usepgfplotslibrary{fillbetween}
\usepackage[np]{numprint}
\usepackage{listings,makecell,pas-tableur,twoopt}
\let\barre\undefined
\usepackage[pstricks]{tablvar}
\definecolor{darkgreen}{RGB}{0,100,0}
\providecommand{\red}{\color{red}}
\providecommand{\blue}{\color{blue}}
\providecommand{\green}{\color{green}}
\providecommand{\black}{\color{black}}
\providecommand{\white}{\color{white}}
\providecommand{\gray}{\color{gray}}
\providecommand{\cyan}{\color{cyan}}
\providecommand{\magenta}{\color{magenta}}
\providecommand{\yellow}{\color{yellow}}
\providecommand{\violet}{\color{violet}}
\providecommand{\orange}{\color{orange}}
\providecommand{\euro}{\texteuro}
\providecommand{\EUR}{\texteuro}
\providecommand{\ieme}{\textsuperscript{e}}
\providecommand{\degres}{\ensuremath{^\circ}}
\providecommand{\degree}{\ensuremath{^\circ}}
\providecommand{\vv}[1]{\overrightarrow{#1}}
\providecommand{\Coord}[1]{#1}
\providecommand{\decosix}{}
\providecommand{\starredbullet}{\ensuremath{\bullet}}
\pagestyle{empty}
\setlength{\parindent}{0pt}
\usepackage{adjustbox}\DeclareUnicodeCharacter{2236}{:}\DeclareUnicodeCharacter{2190}{\ensuremath{\leftarrow}}
\providecommand{\R}{\mathbb{R}}
\providecommand{\N}{\mathbb{N}}
\providecommand{\Z}{\mathbb{Z}}
\providecommand{\Q}{\mathbb{Q}}
\providecommand{\D}{\mathbb{D}}
\providecommand{\C}{\mathbb{C}}
\providecommand{\e}{{\mathrm{e}}}
\providecommand{\vect}[1]{\overrightarrow{#1}}
\providecommand{\vectt}[1]{\overrightarrow{\mathrm{#1}}}
\providecommand{\vvt}[1]{\overrightarrow{\mathrm{#1}}}
\providecommand{\np}[2][]{{#2~#1}}
\providecommand{\vv}[1]{{\overrightarrow{#1}}}
\renewcommand{\d}{{\mathrm{d}}}
\providecommand{\degres}{{^\circ}}
\providecommand{\degre}{{^\circ}}
\providecommand{\blue}{\color{blue}}
\providecommand{\red}{\color{red}}
\providecommand{\enskip}{\quad}
\providecommand{\llbracket}{\lbrack\!\lbrack}
\providecommand{\rrbracket}{\rbrack\!\rbrack}
\providecommand{\euro}{{\text{€}}}
\providecommand{\no}{{\text{n°}}}
\providecommand{\interval}[2]{{[#1;#2]}}
\providecommand{\intervFO}[2]{{[#1;#2[}}
\providecommand{\dotfill}{\ldots}

\providecommand{\ds}{\displaystyle}
\providecommand{\textding}[1]{#1}
\providecommand{\Oij}{\(\mathrm{O};\vec{\imath},\vec{\jmath}\)}
\providecommand{\Oijk}{\(\mathrm{O};\vec{\imath},\vec{\jmath},\vec{k}\)}
\providecommand{\Ouv}{\(\mathrm{O};\vec{u},\vec{v}\)}
\providecommand{\intervFF}[2]{\ensuremath{[#1;#2]}}
\renewcommand{\euro}{\ifmmode\text{\texteuro}\else\texteuro\fi}
\DeclareUnicodeCharacter{00B0}{\ifmmode^{\circ}\else\textdegree\fi}
\DeclareUnicodeCharacter{20AC}{\ifmmode\text{\texteuro}\else\texteuro\fi}
% Énoncés : Éducation nationale. Transcription et corrigés : APMEP, Denis Vergès et contributeurs. Découpage et classement : CoopMaths, Sébastien Lozano, Éric Elter et contributeurs. Adaptation de la présentation : Maths Explorer.
% CC BY-SA : https://coopmaths.fr/www/about/
% https://coopmaths.fr/alea/static/bac/2023/tex/bac_2023_09_sujet2_ameriquesud_4.tex
\begin{document}
\noindent\textit{APMEP / CoopMaths. CC BY-SA. Présentation adaptée par Maths Explorer.}\par\medskip
\medskip

Soit la fonction $f$ définie sur $\R$ par 
\[f(x) = \ln \left(1 + \text{e}^{-x}\right) + \dfrac14 x.\]

On note $\mathcal{C}_f$ la courbe représentative de la fonction $f$ dans un repère orthonormé \Oij{} du plan.

\medskip

\textbf{Partie A}

\medskip

\begin{enumerate}
\item Déterminer la limite de $f$ en $+ \infty$.
\item On admet que la fonction $f$ est dérivable sur $\R$ et on note $f'$ sa fonction dérivée.
	\begin{enumerate}
		\item Montrer que, pour tout réel $x,\: f'(x) = \dfrac{\e^x - 3}{4\left(\e^x + 1\right)}$.
		\item En déduire les variations de la fonction $f$ sur $\R$.
		\item Montrer que l'équation $f(x) = 1$ admet une unique solution $\alpha$ dans l'intervalle [2~;~5].
	\end{enumerate}
\end{enumerate}

\medskip

\textbf{Partie B}

\medskip

On admettra que la fonction $f'$ est dérivable sur $\R$ et pour tout réel $x$,

\[f''(x) = \dfrac{\e^x}{\left(\e^x + 1\right)^2}.\]

On note $\Delta$ la tangente à la courbe $\mathcal{C}_f$ au point d'abscisse $0$.

Dans le graphique ci-dessous, on a représenté la courbe $\mathcal{C}_f$ la tangente $\Delta$ et le quadrilatère MNPQ tel que M et N sont les deux points de la courbe $\mathcal{C}_f$ d'abscisses respectives $\alpha$ et $- \alpha$, et Q et P sont les deux points de la droite $\Delta$ d'abscisses respectives $\alpha$ et $- \alpha$.


\psset{unit=1.2cm,arrowsize=2pt 3}
\begin{center}
\begin{pspicture}(-6,-1)(6,5)
\pspolygon[linewidth=1.25pt,linecolor=blue,fillcolor=lightgray](3.922,-0.3)(3.922,1)(- 3.922,2.961)(-3.922,1.7)

\psaxes[linewidth=1.25pt,Dx=10,Dy=10](0,0)(-6,-1)(6,5)
\psaxes[linewidth=1.25pt,Dx=10,Dy=10]{->}(0,0)(1,1)
\psplot[plotpoints=2000,linewidth=1.25pt,linecolor=red]{-6}{6}{1 1 2.71828 x exp div add ln 0.25 x mul add}
\psplotTangent[linewidth=1.25pt,linecolor=blue]{0}{6}{1 1 2.71828 x exp div add ln 0.25 x mul add}

\psline[linestyle=dashed](-3.922,1.7)(-3.922,0)
\uput[d](0.5,0){$\vect{\imath}$}\uput[l](0,0.5){$\vect{\jmath}$}
\uput[d](3.922,-0.3){Q}\uput[u](3.922,1){M}\uput[ur](- 3.922,2.961){N}\uput[dl](-3.922,1.7){P}
\uput[dr](3.922,0){$\alpha$}\uput[d](-3.922,0){$- \alpha$}\uput[dl](0,0){O}\uput[d](-5.2,2){\blue$\Delta$}
\uput[d](-5.2,3.8){\red $\mathcal{C}_f$}
\end{pspicture}
\end{center}

\smallskip

\begin{enumerate}
\item 
	\begin{enumerate}
		\item Justifier le signe de $f''(x)$ pour $x \in \R$.
		\item En déduire que la portion de la courbe $\mathcal{C}_f$ sur l'intervalle $[-\alpha~;~\alpha]$, est inscrite dans le quadrilatère MNPQ.
	\end{enumerate}
\item 	
	\begin{enumerate}
		\item Montrer que $f(-\alpha) = \ln \left(\e^{-\alpha} + 1\right) + \dfrac34 \alpha$.
		\item Démontrer que le quadrilatère MNPQ est un parallélogramme.
	\end{enumerate}
\end{enumerate}

\end{document}
